% !TeX root = ../queens.tex
\section{Upper queens and sorted lower rows}\label{sec:rows}
Write $(x_k^U,y_k^U)$ and $(x_k^L,y_k^L)$ for the coordinates of the $k$th
upper and lower queens respectively, whenever they exist. Each family is
ordered by increasing column, starting at $k=1$.
For $k\ge1$, we call $\{(x,y)\in\N^2:y-x=k\}$ and
$\{(x,y)\in\N^2:y-x=-k\}$ respectively the \emph{$k$th upper}
and \emph{$k$th lower diagonal}; $\{(x,y)\in\N^2:x+y=k\}$
is the \emph{$k$th antidiagonal} (Figure~\ref{fig:infinite-board}).
A row or column is \emph{upper} if it contains an upper queen, and \emph{lower} if it contains a lower queen.

The next identity is implicit in Knuth's program \texttt{infty-queens}
\cite{KnuthQueens}. 
\begin{lemma}\label{lem:upper}
The $k$th upper queen lies on the $k$th upper diagonal; that is,
$y_k^U=x_k^U+k$.
\end{lemma}
\noindent
\begin{minipage}[t]{0.60\linewidth}
\vspace{0pt}
\begin{proof}
Suppose inductively that the first $k-1$ upper queens have used
the upper diagonals $1,\dots,k-1$ respectively, and consider column $n=x_k^U$. Those upper diagonals prevent rows $n+1,\dots,n+k-1$ from being used,
so the lowest upper square that is possibly free is
$(n,n+k)$. By induction, the $k$th upper diagonal that it lies on is free,
so it remains
to show that its row and antidiagonal are free as well. 
Every earlier upper queen $(x_j^U,y_j^U)$ has $j<k$ and
$x_j^U<n$, so the inductive hypothesis gives
$y_j^U=x_j^U+j<n+k$. Every earlier lower queen, and the queen at the
origin, also lies below row $n+k$. Thus every earlier queen has both
coordinates strictly smaller than those of $(n,n+k)$, so none shares
its row or antidiagonal. Since the queen in column $n$ is upper and
all smaller upper rows are attacked, the greedy rule chooses row
$n+k$. This completes the induction.
\end{proof}
\end{minipage}\hfill
\begin{minipage}[t]{0.38\linewidth}
\vspace{-0.5\baselineskip}
\centering
\setboardfontsize{11pt}
\begin{tikzpicture}[x=0.46cm,y=0.46cm,font=\scriptsize,line cap=round,
  earlier queen/.style={inner sep=0pt,fill=white,text=queensUpper},
  free line/.style={dash pattern=on 3pt off 2pt,line width=0.5pt}]
  % Schematic of the step n=7, k=5, with symbolic labels.
  \draw[black!35,line width=0.5pt] (0,0) -- (7,7);
  \draw[black!55,line width=0.5pt] (7,-0.1) -- (7,12);
  \node[anchor=north] at (7,-0.3) {$n$};
  % Previously occupied upper diagonals meet column n in forbidden rows.
  \foreach \j in {1,...,4} {
    \draw[queensUpper!65,line width=0.6pt] (0,\j) -- (7,{7+\j});
    \node[text=black!55,inner sep=0pt] at (7,{7+\j}) {$\times$};
  }
  % The candidate's three attack lines are all unoccupied.
  \draw[free line,queensUpper] (0,5) -- (9,14);
  \node[rotate=45,anchor=south,text=queensUpper,inner sep=2pt]
    at (2.8,7.8) {upper diagonal $k$};
  \draw[free line,black!40] (0,12) -- (10.8,12);
  \node[anchor=south west,inner sep=2pt,yshift=-0.8pt] at (7.5,12.05) {row $n+k$};
  \draw[free line,black!45] (4,15) -- (11.8,7.2);
  \node[anchor=west,inner sep=2pt,yshift=0.8pt] at (7.3,8) {$n+1$};
  \node[inner sep=0pt] at (8.6,9.5) {$\vdots$};
  \node[anchor=west,inner sep=2pt,yshift=0.8pt] at (7.3,11) {$n+k-1$};
  \node[align=center,anchor=north,inner sep=1pt] at (10.5,6.9)
    {antidiagonal\\$2n+k$};
  % All earlier queens lie strictly below and to the left of the candidate.
  \foreach \x/\y in {1/2,2/4,5/8,6/10} {
    \node[earlier queen] at (\x,\y) {\BlackQueenOnWhite};
  }
  \foreach \x/\y in {3/1,4/3} {
    \node[earlier queen,text=queensLower] at (\x,\y) {\BlackQueenOnWhite};
  }
  \node[earlier queen,text=black] at (0,0) {\BlackQueenOnWhite};
  \filldraw[fill=white,draw=queensUpper,line width=0.8pt]
    (6.8,11.8) rectangle (7.2,12.2);
  \node[font=\footnotesize,anchor=south,yshift=-4.5pt] at (6,15.3) {$(n,n+k)$};
  \draw[->,>=stealth,black!55,line width=0.4pt]
    (6.3,15) .. controls (6.6,14.3) and (6.6,12.9) .. (6.95,12.25);
\end{tikzpicture}
\end{minipage}
\par\medskip
Note that Lemma~\ref{lem:upper} implies that successive upper rows differ by at least two.

The permutation property below follows from the row-permutation
argument of Rob Pratt, Bob Selcoe, and N.~J.~A.~Sloane (July~2, 2016)
in OEIS entry~\href{https://oeis.org/A269526}{A269526} \cite[Theorem~R]{OEIS}, reproduced in
\cite[Theorem~23]{DSS}; see also Knuth
\cite[exercise~7.2.2.1--37]{Knuth4B}. 
In combinatorial game theory terms, that theorem says that every
nonnegative Sprague--Grundy value occurs exactly once in each row and
column. The lemma below is its zero-value case. 
\begin{lemma}\label{lem:surjective}
The map $n\mapsto q_n$ is a permutation of $\N$.
\end{lemma}
\begin{proof}
The rows $q_n$ are distinct by construction, so it suffices to show that
every row is occupied by some queen. Suppose otherwise, and let $m$ be the least unoccupied row. 
For $n>m$, an earlier queen $(i,q_i)$ can attack $(n,m)$ diagonally
only if
\[
 n-m=i-q_i,\qquad q_i=m+i-n<m.
\]
Since there are at most $m$ such earlier queens, it follows that $(n,m)$ has no diagonal attackers for sufficiently large $n$. 
Now after some finite time, every row below $m$ has been occupied.
Since row $m$ remains unused by hypothesis,
the greedy rule requires $n+m$ to be an occupied antidiagonal for sufficiently large $n$. Hence the
set of occupied antidiagonals contains a tail $[N\dts\infty)$ of the integers.

We now show that such a tail is impossible. Choose an integer $t>2N+1$
and consider the $t$ queens occupying the antidiagonals
$N,N+1,\ldots,N+t-1$. These queens lie in distinct columns and distinct
rows. If we arrange their column coordinates in increasing order, the
first is at least $0$, the second at least $1$, and so on.

\par\medskip\noindent
\begin{minipage}[t]{0.38\linewidth}
\vspace{0pt}
\centering
\setboardfontsize{12pt}
\begin{tikzpicture}[x=0.43cm,y=0.43cm,font=\scriptsize,line cap=round,
  coordinate guide/.style={black!30,dash pattern=on 2pt off 2pt,line width=0.4pt},
  count queen/.style={inner sep=0pt,text=black}]
  % A possible finite block, not an instance of the contradictory hypothesis.
  % The axis labels are rank bounds, not the drawn coordinate values.
  \fill[queensUpper!7] (0,6) -- (0,9) -- (9,0) -- (6,0) -- cycle;
  \foreach \s in {7,8} {
    \draw[queensUpper!35,line width=0.45pt] (0,\s) -- (\s,0);
  }
  \foreach \s in {6,9} {
    \draw[queensUpper!75,line width=0.65pt] (0,\s) -- (\s,0);
  }
  \foreach \x/\y in {1/5,3/4,6/2,9/0} {
    \draw[coordinate guide] (\x,0) -- (\x,\y) -- (0,\y);
  }
  \draw[black!65,line width=0.5pt] (0,9.6) -- (0,0) -- (9.6,0);
  \node[anchor=south] at (0,9.65) {$y$};
  \node[anchor=west] at (9.65,0) {$x$};
  \foreach \x/\bound in {1/{\ge0},3/{\ge1},6/{\cdots},9/{\ge t-1}} {
    \draw[black!65,line width=0.4pt] (\x,0) -- (\x,-0.15);
    \node[anchor=north,inner sep=2pt] at (\x,-0.22) {$\bound$};
  }
  \foreach \y/\bound in {0/{\ge0},2/{\ge1},4/{\vdots},5/{\ge t-1}} {
    \draw[black!65,line width=0.4pt] (0,\y) -- (-0.15,\y);
    \node[anchor=east,inner sep=2pt] at (-0.22,\y) {$\bound$};
  }
  \node[text=queensUpper,rotate=-45,anchor=south,inner sep=2pt]
    at (3.6,2.4) {$N$};
  \node[text=queensUpper,rotate=-45,anchor=south,inner sep=2pt]
    at (4.5,4.5) {$N+t-1$};
  \foreach \x/\y in {1/5,3/4,6/2,9/0} {
    \node[count queen] at (\x,\y) {\BlackQueenOnWhite};
  }
\end{tikzpicture}
%\par\smallskip
{\scriptsize One queen per antidiagonal (schematic).\par
The axis labels give lower bounds\par
on the coordinates in increasing order.\par}
\end{minipage}\hfill
\begin{minipage}[t]{0.60\linewidth}
\vspace{0pt}
Thus the sum of their column coordinates is at least
\[
 0+1+\cdots+(t-1)=\frac{t(t-1)}2.
\]
The same bound holds for their row coordinates. Thus the sum of all
their row and column coordinates is at least $t(t-1)$.

On the other hand, each queen's row coordinate plus its column coordinate
is precisely the label of its antidiagonal. Summing these labels gives
the same total as
\[
 N+(N+1)+\cdots+(N+t-1)=tN+\frac{t(t-1)}2.
\]
\end{minipage}
\par\medskip
\noindent Consequently,
\[
 t(t-1)\le tN+\frac{t(t-1)}2,
 \qquad\text{so}\qquad t\le2N+1.
\]
This contradicts our choice of $t$. Hence every row is occupied, as required.
\end{proof}

For $n\in\N$, define
\begin{equation}\label{eq:counts}
 U(n)=\#\{1\le i\le n:q_i>i\},\qquad L(n)=n-U(n).
\end{equation}
These count upper and lower queens through column $n$. There are infinitely many lower queens: successive upper rows differ by at least two by Lemma~\ref{lem:upper}, omitting infinitely many positive integers, and every row is occupied by Lemma~\ref{lem:surjective}.
Their columns satisfy $x_1^L<x_2^L<\cdots$, though the row values $y_j^L=q_{x_j^L}$ need not be increasing. We define the \emph{lower-diagonal magnitudes}
\begin{equation}\label{eq:rank-definition}
 d_j\coloneqq x_j^L-y_j^L>0\qquad(j\ge1),
\end{equation}
so that the $j$th lower queen lies on the $d_j$th lower diagonal.
The $d_j$ are distinct, and they are not increasing. Corollary~\ref{cor:diagonal-coverage}
will show that they form a permutation of the positive integers.

Let $v_1<v_2<\cdots$ be the lower-row values $y_j^L$, sorted into increasing order. We call $v_j$ the $j$th \emph{sorted} lower row
and $y_j^L$ the $j$th \emph{chronological} lower row (or just $j$th lower row).


\begin{center}
\begin{minipage}{\linewidth}
\centering
\renewcommand{\arraystretch}{1.15}
\setlength{\tabcolsep}{6pt}
\begin{tabular}{@{}l*{15}{r}@{}}
\toprule
$j$     & 1 & 2 & 3 & 4 & 5 & 6 & 7 & 8 & 9 & 10 & 11 & 12 & 13 & 14 & 15 \\
\midrule
$x_j^L$ & 3 & 4 & 9 & 10 & 12 & 14 & 19 & 20 & 25 & 26 & 28 & 30 & 33 & 37 & 39 \\
$y_j^L$ & 1 & 3 & 5 & 7 & 6 & 9 & 11 & 13 & 15 & 17 & 16 & 19 & 20 & 22 & 23 \\
$d_j$   & 2 & 1 & 4 & 3 & 6 & 5 & 8 & 7 & 10 & 9 & 12 & 11 & 13 & 15 & 16 \\
$v_j$   & 1 & 3 & 5 & 6 & 7 & 9 & 11 & 13 & 15 & 16 & 17 & 19 & 20 & 22 & 23 \\
$U(j)$  & 1 & 2 & 2 & 2 & 3 & 4 & 5 & 6 & 6 & 6 & 7 & 7 & 8 & 8 & 9 \\
\bottomrule
\end{tabular}
\par
\captionof{table}{Lower-queen sequences and upper counts for $1\le j\le15$.}
\label{tab:lower-values}
\end{minipage}
\end{center}
\vspace{\baselineskip}

% If the illustrated proof moves to the next page, leave space at the foot
% of this page instead of stretching the preceding displayed equations.
% \par\vfil\penalty-200\vfilneg
\begin{samepage}
The following identity is a special case of the Lambek--Moser theorem on
inverse and complementary sequences \cite{LM}, applied to $f(k)=x_k^U$. 
\begin{lemma}\label{lem:complement}
For every $j\ge1$, the $j$th sorted lower row $v_j$ satisfies
\begin{equation}\label{eq:sorted-rows}
 v_j=j+U(j-1).
\end{equation}
\end{lemma}
\noindent
\begin{minipage}[t]{0.60\linewidth}
\vspace{0pt}
\begin{proof}
Put $h=U(j-1)$ and $r=j+h$. We will show that row $r$ is a lower row
and that exactly $j$ lower rows lie in $[1\dts r]$. This identifies $r$
as the $j$th smallest lower row, namely $v_j$.

There are exactly $h$ upper queens in columns to the left of $j$. Each has an
index $k\le h$ and a column $x_k^U\le j-1$, so Lemma~\ref{lem:upper} gives
\[
 y_k^U=x_k^U+k\le(j-1)+h=r-1.
\]
Every upper queen in column $j$ or later instead has $k\ge h+1$ and
$x_k^U\ge j$, so
\[
 y_k^U=x_k^U+k\ge j+(h+1)=r+1.
\]
Thus no upper queen occupies row $r$, and exactly $h$ of the rows in
$[1\dts r]$ are upper rows.

By Lemma~\ref{lem:surjective}, every one of these positive rows is
occupied. The remaining $r-h=j$ rows are therefore lower rows, and
row $r$ itself is one of them. Hence $r$ is the largest of these $j$
lower rows, proving $v_j=r=j+U(j-1)$.
\end{proof}
\end{minipage}\hfill
\begin{minipage}[t]{0.38\linewidth}
\vspace{-0.5\baselineskip}
\centering
\setboardfontsize{11pt}
\begin{tikzpicture}[x=0.34cm,y=0.34cm,font=\scriptsize,line cap=round,
  example queen/.style={inner sep=0pt,fill=white},
  example guide/.style={dash pattern=on 3pt off 2pt,line width=0.5pt}]
  % Actual queen positions in columns 0 through 12, with j=5 and r=7.
  \node[font=\footnotesize,anchor=south] at (5,19.2) {$h=U(4)=2$};
  \foreach \y in {1,3,5,6} {
    \draw[queensLower!22,line width=0.35pt] (0,\y) -- (12.6,\y);
  }
  \foreach \y in {2,4} {
    \draw[queensUpper!30,line width=0.35pt] (0,\y) -- (12.6,\y);
  }
  \draw[black!55,line width=0.5pt] (0,19) -- (0,0) -- (13,0);
  \node[anchor=south] at (0,19.1) {$y$};
  \node[anchor=west] at (13,0) {$x$};
  \foreach \y in {1,3,5,6,7} {
    \node[anchor=east,text=queensLower,inner sep=2pt] at (-0.2,\y) {$\y$};
  }
  \foreach \y in {2,4,8,10,12,14,18} {
    \node[anchor=east,text=queensUpper,inner sep=2pt] at (-0.2,\y) {$\y$};
  }
  \foreach \x in {10,12} {
    \draw[black!55,line width=0.4pt] (\x,0) -- (\x,-0.2);
    \node[anchor=north,inner sep=2pt] at (\x,-0.25) {$\x$};
  }
  \draw[example guide,black!45] (5,0) -- (5,19);
  \node[anchor=north,inner sep=2pt] at (5,-0.25) {$j=5$};
  \draw[example guide,queensLower,line width=0.65pt] (0,7) -- (12.8,7);
  \node[anchor=west,text=queensLower,inner sep=2pt,yshift=0.8pt] at (12.8,7) {$r=7$};
  \foreach \x/\y in {1/2,2/4,5/8,6/10,7/12,8/14,11/18} {
    \node[example queen,text=queensUpper] at (\x,\y) {\BlackQueenOnWhite};
  }
  \foreach \x/\y in {3/1,4/3,9/5,12/6,10/7} {
    \node[example queen,text=queensLower] at (\x,\y) {\BlackQueenOnWhite};
  }
  \node[example queen,text=black] at (0,0) {\BlackQueenOnWhite};
  % Keep the diagonal continuous over the queen nodes' white backgrounds.
  \draw[black!35,line width=0.5pt] (0,0) -- (12.6,12.6);
  \node[anchor=north,align=center,inner sep=2pt] at (6.5,-2.3)
    {Upper rows through $7$: \textcolor{queensUpper}{$2,4$}.\\
     Lower rows through $7$: \textcolor{queensLower}{$1,3,5,6,7$}.\\[2pt]
     Thus $v_5=7=5+U(4)$.};
\end{tikzpicture}
\end{minipage}
\par\medskip
\end{samepage}

